%% kapa: exam=1 round=35 year=2024 period=1 form=A subject=economics subject_ko=경제학원론 seq=26 number=66
%% source: Q-Net 감정평가사 기출문제 (cst00309) artlSeq=5213229 file=raw/1008/5213229/1교시 문제지 원본.pdf
%% conversion: kordoc 4.24.51 -> tools/gen1.py (auto); review=proofread-sonnet; flags=none
\begin{question}{66}
통화정책에 관한 설명으로 옳지 않은 것은?
\begin{choices}
  \item 공개시장 매입은 본원통화를 증가시켜 이자율을 하락시킨다.
  \item 재할인율 인상은 재할인대출을 감소시켜 이자율을 상승시킨다.
  \item 자산가격경로는 이자율이 하락할 경우 자산가격이 상승하여 부(富)의 효과로 소비가 증가하는 경로이다.
  \item 신용경로는 중앙은행이 화폐공급을 축소할 경우 은행대출이 감소되어 기업투자와 가계소비가 위축되는 경로이다.
  \item 환율경로는 이자율이 상승할 경우 자국통화가치가 하락하여 순수출이 증가하는 경로이다.
\end{choices}
\end{question}
\answer{5}
